%
\subsection*{Estados y sistemas equivalentes}
\label{CH07-B01}
Una EDO de segundo orden $y''=f(t,y,y')$ se escribe como
$x_1'=x_2$, $x_2'=f(t,x_1,x_2)$, con $x_1=y$, $x_2=y'$ y ambos
datos iniciales. El producto matriz-vector se calcula por filas:
$(A\mathbf x)_i=\sum_j a_{ij}x_j$.

Eliminar una componente puede producir una EDO escalar, pero la
equivalencia exige reconstruirla con las mismas constantes y comprobar
todo el sistema. Las condiciones iniciales de la ecuación eliminada
proceden del sistema original. Si éste tiene forma
$B\mathbf x'=A\mathbf x+\mathbf f$, sólo puede despejarse
$\mathbf x'$ mediante $B^{-1}$ cuando $B$ es invertible. Una matriz
$B$ singular puede imponer una relación algebraica entre estados y datos.
Recuérdense sistemas lineales y determinantes.
Consulta: solución de \exref{M2-E0238}.

%
\subsection*{Evolución lineal y autovectores}
\label{CH07-B02}
En $\mathbf x'=A\mathbf x$, con $A$ constante, un autovector no nulo
$\mathbf v$ de autovalor $\lambda$ produce
$e^{\lambda t}\mathbf v$. Calcúlense los autovalores con
$\det(A-\lambda I)=0$ y los autovectores resolviendo
$(A-\lambda I)\mathbf v=0$. Para una matriz real, las partes real e
imaginaria de un modo complejo conjugado dan soluciones reales.

Una raíz repetida no garantiza suficientes autovectores independientes.
En el caso concreto $A=\lambda I+N$, $N^2=0$, se puede comprobar por
derivación que
\[
 \mathbf x(t)=e^{\lambda(t-t_0)}
 [I+(t-t_0)N]\mathbf x(t_0).
\]
Para $\mathbf x'=A\mathbf x+\mathbf f(t)$, la solución es homogénea
más particular. Una entrada constante sugiere una particular constante
sólo si $A\mathbf x_p+\mathbf f=0$ es compatible; otras entradas
pueden requerir otro ensayo. Verifíquense todas las componentes y datos.
Consultas: soluciones de \exref{M2-E0047} y \exref{M2-E0065}.

%
\subsection*{Balances, reducción y observación}
\label{CH07-B03}
Escríbase un balance por cada cantidad almacenada. En una mezcla,
$Q_i=V_i c_i$; si cambia $V_i$, también interviene $V_i'c_i$.
Al sumar balances, los intercambios internos se cancelan y quedan los
intercambios con el exterior y las fuentes o consumos. Que aumente una
componente no implica que aumente el total.

Las variables suma y diferencia pueden separar evolución conjunta y
desequilibrio. Para justificar que dos estados permanecen iguales,
compruébese que sus derivadas coinciden cuando los estados coinciden;
con datos iguales y unicidad, esa igualdad se conserva. No basta imponerla
en una sola ecuación. Distínganse el estado físico, la salida de un sensor
$H(\mathbf x)$ y un registro acumulado
$z'=H(\mathbf x)$: miden objetos distintos y pueden tener unidades
distintas. Consulta: solución de \exref{M2-E0343}.

%
\subsection*{Etapas vectoriales y límites de la aproximación}
\label{CH07-B04}
Euler, Heun y RK4 se aplican a $\mathbf x'=\mathbf f(t,\mathbf x)$
con las mismas fórmulas del capítulo 3, sustituyendo números por vectores.
Cada etapa utiliza todas las componentes del mismo estado provisional;
actualizar una componente y usarla de inmediato para calcular otra cambia
el método.

Compruébense invariantes, signos, unidades y balances. Conservar una suma
no garantiza aproximar bien cada componente. Las primeras etapas manuales
permiten entender el algoritmo; cuando se pide una malla larga, puede usarse
la ayuda de cálculo de esta edición para las repeticiones.
La concordancia entre mallas es evidencia numérica, no una tolerancia
certificada ni una validación empírica de las leyes. Para una decisión
garantizada, búsquense cotas justificadas mediante balances y
desigualdades, e indíquense las hipótesis del modelo.
Consultas: soluciones de \exref{M2-E0315} y \exref{M2-E0362}.
